% problem_id: S001-theorem-quasi-DM
% source_id: S001 (Stacks Project)
% source_file: stacks-morphisms.tex
% statement_locator: stacks-morphisms.tex lines 4072--4080
% proof_locator: stacks-morphisms.tex lines 4082--4190
% env: theorem
% id_note: label 在本来源内唯一
% status: complete (statement + author proof verbatim + reason + locators)
%
\begin{theorem}
\label{theorem-quasi-DM}
Let $\mathcal{X}$ be an algebraic stack. The following are equivalent
\begin{enumerate}
\item $\mathcal{X}$ is quasi-DM, and
\item there exists a scheme $W$ and a surjective, flat, locally finitely
presented, locally quasi-finite morphism $W \to \mathcal{X}$.
\end{enumerate}
\end{theorem}

\begin{proof}
The implication (2) $\Rightarrow$ (1) is
Lemma \ref{lemma-properties-covering-imply-diagonal}.
Assume (1).
Let $x \in |\mathcal{X}|$ be a finite type point. We will produce a scheme
over $\mathcal{X}$ which ``works'' in a neighbourhood of $x$. At the end
of the proof we will take the disjoint union of all of these to conclude.

\medskip\noindent
Let $U$ be an affine scheme, $U \to \mathcal{X}$ a smooth morphism, and
$u \in U$ a closed point which maps to $x$, see
Lemma \ref{lemma-point-finite-type}.
Denote $u = \Spec(\kappa(u))$ as usual. Consider the following
commutative diagram
$$
\xymatrix{
u \ar[d] & R \ar[l] \ar[d] \\
U \ar[d] & F \ar[d] \ar[l]^p \\
\mathcal{X} & u \ar[l]
}
$$
with both squares fibre product squares, in particular
$R = u \times_\mathcal{X} u$. In the proof of
Lemma \ref{lemma-point-finite-type-monomorphism}
we have seen that $(u, R, s, t, c)$ is a groupoid in algebraic spaces
with $s, t$ locally of finite type. Let $G \to u$ be the stabilizer group
algebraic space (see
Groupoids in Spaces, Definition
\ref{spaces-groupoids-definition-stabilizer-groupoid}).
Note that
$$
G = R \times_{(u \times u)} u =
(u \times_\mathcal{X} u) \times_{(u \times u)} u =
\mathcal{X} \times_{\mathcal{X} \times \mathcal{X}} u.
$$
As $\mathcal{X}$ is quasi-DM we see that
$G$ is locally quasi-finite over $u$. By
More on Groupoids in Spaces, Lemma
\ref{spaces-more-groupoids-lemma-groupoid-on-field-dimension-equal-stabilizer}
we have $\dim(R) = 0$.

\medskip\noindent
Let $e : u \to R$ be the identity of the groupoid. The composition
$u \to R \to u$ is equal to the identity morphism of $u$.
Note that $R \subset F$ is a closed
subspace as $u \subset U$ is a closed subscheme. Hence we can also think
of $e$ as a point of $F$; call it $z$.
Consider the maps of \'etale local rings
$$
\mathcal{O}_{U, u}
\xrightarrow{p^\sharp}
\mathcal{O}_{F, \overline{z}}
\longrightarrow
\mathcal{O}_{R, \overline{e}}
$$
Note that $\mathcal{O}_{R, \overline{e}}$ has dimension $0$ by the result
of the first paragraph. On the other hand, the kernel of the second arrow is
$p^\sharp(\mathfrak m_u)\mathcal{O}_{F, \overline{z}}$ as
$R$ is cut out in $F$ by $\mathfrak m_u$. Thus we see that
$$
\mathfrak m_{\overline{z}} =
\sqrt{p^\sharp(\mathfrak m_u)\mathcal{O}_{F, \overline{z}}}
$$
On the other hand, as the morphism $U \to \mathcal{X}$ is smooth
we see that $F \to u$ is a smooth morphism of algebraic spaces.
This means that $F$ is a regular algebraic space
(Spaces over Fields, Lemma \ref{spaces-over-fields-lemma-smooth-regular}).
Hence $\mathcal{O}_{F, \overline{z}}$ is a regular local ring
(Properties of Spaces, Lemma \ref{spaces-properties-lemma-regular}).
Note that a regular local ring is Cohen-Macaulay
(Algebra, Lemma \ref{algebra-lemma-regular-ring-CM}).
Let $d = \dim(\mathcal{O}_{F, \overline{z}})$. By
Algebra, Lemma \ref{algebra-lemma-find-sequence-image-regular}
we can find $f_1, \ldots, f_d \in \mathcal{O}_{U, u}$ whose images
$\varphi(f_1), \ldots, \varphi(f_d)$ form a regular sequence
in $\mathcal{O}_{F, \overline{z}}$. By
Lemma \ref{lemma-slice}
after shrinking $U$ we may assume that
$Z = V(f_1, \ldots, f_d) \to \mathcal{X}$ is flat and
locally of finite presentation. Note that by construction
$F_Z = Z \times_\mathcal{X} u$ is a closed subspace of
$F = U \times_\mathcal{X} u$, that $z$ is a point of this closed subspace,
and that
$$
\dim(\mathcal{O}_{F_Z, \overline{z}}) = 0.
$$
By
Morphisms of Spaces,
Lemma \ref{spaces-morphisms-lemma-dimension-fibre-at-a-point}
it follows that $\dim_z(F_Z) = 0$ because the transcendence degree
of $z$ relative to $u$ is zero. Hence it follows from
Lemma \ref{lemma-quasi-finite-at-point}
that after possibly shrinking $U$ the morphism $Z \to \mathcal{X}$
is locally quasi-finite.

\medskip\noindent
We conclude that for every finite type point $x$ of $\mathcal{X}$ there
exists a locally quasi-finite, flat, locally finitely presented
morphism $f_x : Z_x \to \mathcal{X}$ with $x$ in the image of $|f_x|$.
Set $W = \coprod_x Z_x$ and $f = \coprod f_x$. Then $f$ is flat, locally
of finite presentation, and locally quasi-finite. In particular the
image of $|f|$ is open, see
Properties of Stacks, Lemma \ref{stacks-properties-lemma-topology-points}.
By construction the image contains all finite type points of $\mathcal{X}$,
hence $f$ is surjective by
Lemma \ref{lemma-enough-finite-type-points} (and
Properties of Stacks, Lemma
\ref{stacks-properties-lemma-characterize-surjective}).
\end{proof}
